\documentclass[11pt]{article}
\usepackage[margin=1in]{geometry}
\usepackage{times}
\usepackage{setspace}
\setstretch{1.15}
\usepackage{graphicx}
\usepackage{booktabs}
\usepackage{amsmath}
\usepackage{amssymb}
\usepackage{caption}
\usepackage{xcolor}
\usepackage[colorlinks=true, linkcolor=blue!70!black, citecolor=blue!70!black, urlcolor=blue!70!black]{hyperref}
\urlstyle{same}
\renewcommand{\ttdefault}{cmtt}

\def\labelitemi{--}

\begin{document}

\begin{center}
    {\Large\bfseries Online Appendices: Opening the Black Box of Large Language Models in Text-based Measurement\par}
    \vspace{0.5em}
    {\large\itshape Staged Pipelines with Verifiable Checkpoints\par}
    \vspace{1em}
    {\normalsize\bfseries Milan Miric, Albert Roh, Nan Jia\par}
    \vspace{0.2em}
    {\small\itshape University of Southern California, Marshall School of Business\par}
    \vspace{0.8em}
    {\footnotesize\textbf{Proposal for Special Issue:} AI-Enabled Frontiers in Organizational Science (\textit{Organization Science})\par}
    \vspace{0.3em}
    {\footnotesize\textbf{Repository:} \url{https://github.com/HC-Geopolitics-Innovation/LLM-Text-Analysis-and-Synthesis}\par}
\end{center}

\vspace{1.5em}

\section*{Appendix A. Overview of Steps in Large Language Model Enabled Research Pipeline}
\addcontentsline{toc}{section}{Appendix A. Overview of Steps in Large Language Model Enabled Research Pipeline}

\subsection*{A.1. Long-Document Text Summarization}
\subsubsection*{A.1.1. Overview of Processing Pipeline}

Processing long text documents is a challenge for large language models because they are prone to degradation in attention as input text increases (Liu et al., 2023), struggle to preserve factual detail across multiple conversational turns, and can introduce hallucinations during complex open-ended synthesis. 

To overcome these constraints, we propose a staged Map-Reduce summarization pipeline that preserves document integrity:
\begin{enumerate}
    \item \textbf{Data Ingestion \& Speaker Turn Partitioning:} Documents are ingested from source XML/HTML formats, parsed into turn-level records, and attributed to specific speakers.
    \item \textbf{Chunk-Level Map Summarization:} Documents exceeding the model's high-attention context threshold are divided into sequential semantic chunks (roughly 4,000--8,000 tokens) with overlapping boundaries to prevent conversational clipping. Prompts focus the model on identifying substantive factual points.
    \item \textbf{Hierarchical Reduce Synthesis:} Intermediate chunk summaries are aggregated through a structured reduce prompt that reconciles cross-chunk themes into an entity-attributed document summary.
    \item \textbf{Discriminant Validity Checkpoint:} Summaries are validated against external human benchmarks (such as professional C-SPAN editorial abstracts).
\end{enumerate}

\subsubsection*{A.1.2. Application: Summarizing Long Text Congressional Hearings}

We apply this architecture to 2,352 congressional committee hearings before the House Science, Space, and Technology Committee and the Senate Commerce, Science, and Transportation Committee. Transcripts average roughly 40,000 words. Figure 1 reports the cosine similarity between our generated summaries and independent, professional C-SPAN abstracts across same-hearing, same-committee, and different-committee pairs.

\begin{figure}[htbp]
    \centering
    \includegraphics[width=0.85\textwidth]{figures/fig1_cosine_similarity.png}
    \caption{\textbf{Cosine Similarity Between Hand-Coded Summaries and Large Language Model Summaries across Committees.}}
    \vspace{4pt}
    {\footnotesize\textit{Notes: Comparison of cosine similarity distributions between model-generated summaries and independent human-authored C-SPAN abstracts. Same-document pairs exhibit high semantic alignment, whereas different-committee pairs show distinct separation, verifying discriminant validity.}\par}
\end{figure}

\subsubsection*{A.1.3. Research Outputs}

\begin{itemize}
    \item \textbf{Cleaned Transcript Segments:} Fully segmented, speaker-attributed transcripts with timestamp and witness provenance.
    \item \textbf{Intermediate Chunk Summaries:} Auditable, passage-level summaries supporting localized verification.
    \item \textbf{Final Hearing Summaries:} Entity-attributed committee hearing summaries for downstream classification.
\end{itemize}

\subsection*{A.2. Extracting Topics, Positions, and Arguments from Long-Form Documents}
\subsubsection*{A.2.1. Overview of Processing Pipeline}

\begin{enumerate}
    \item \textbf{Taxonomic Clustering \& Codebook Generation:} We use LLMs to extract emerging policy topics from statements and crosswalk them to standardized classification systems like the Comparative Agendas Project (CAP).
    \item \textbf{Stance Extraction:} Identifying discrete positions (support, concern, neutral) held by individual speakers on focal topics.
    \item \textbf{Validation Against External Criteria:} Agreement audits across prompt variants and cross-validation against legislative roll-call voting records.
\end{enumerate}

\subsubsection*{A.2.2. Application: Extracting Stances and Arguments around Focal Topics}

Figure 2 illustrates the distribution of identified congressional stances across major technological domains, highlighting the balance of support versus concern.

\begin{figure}[htbp]
    \centering
    \includegraphics[width=0.90\textwidth]{figures/fig2_topics_stances_4panel.png}
    \caption{\textbf{Topics and Overview of Stance (Support / Concern of Policymakers).}}
    \vspace{4pt}
    {\footnotesize\textit{Notes: Four-panel overview tracking congressional topic prevalence, stance distributions, member polarization, and longitudinal legislative attention across science and technology hearings.}\par}
\end{figure}

\subsubsection*{A.2.3. Research Outputs}

\begin{itemize}
    \item \textbf{Standardized 111-Topic Policy Codebook:} Fully crosswalked to CAP categories.
    \item \textbf{Member-Topic Stance Panel:} Micro-level dataset tracking individual member arguments and policy positions over time.
\end{itemize}

\subsection*{A.3. Policy Relevance for Publicly Listed Firms}
\subsubsection*{A.3.1. Overview of Processing Pipeline}

\begin{enumerate}
    \item \textbf{Corporate Text Extraction:} We extract Item 1 (Business Description) from annual 10-K filings.
    \item \textbf{Centroid Embedding Estimation:} Generating high-dimensional dense embeddings of policy debate stances and firm business summaries.
    \item \textbf{Cosine Distance Scoring:} Measuring semantic cosine similarity between each firm's business profile and the policy centroid, demeaned relative to sample benchmarks.
\end{enumerate}

\subsubsection*{A.3.2. Application: Connecting Focal Concerns and Opportunities of Policymakers to Business Areas of Publicly Listed Firms}

Figure 3 reports topic relevance scores for publicly listed firms regarding Consumer Data Privacy Legislation (Topic 68). Panel A reports industry-level exposure by NAICS sector relative to the sample mean. Panel B reports firm-level relevance scores for selected high- and low-relevance firms.

\begin{figure}[htbp]
    \centering
    \begin{minipage}{0.48\textwidth}
        \centering
        \textbf{\footnotesize Panel A: NAICS Sector Cosine Relative to Mean}\par\vspace{2pt}
        \includegraphics[width=\textwidth]{figures/fig3_panel_a_naics_relevance.png}
    \end{minipage}\hfill
    \begin{minipage}{0.48\textwidth}
        \centering
        \textbf{\footnotesize Panel B: Selected High and Low Relevance Firms}\par\vspace{2pt}
        \includegraphics[width=\textwidth]{figures/fig3_panel_b_top_firms.png}
    \end{minipage}
    \caption{\textbf{Topic Relevance Scores for Publicly Listed Firms.}}
    \vspace{4pt}
    {\footnotesize\textit{Notes: Panel A reports firm summary cosine similarity by NAICS sector relative to the sample mean (Topic 68 policy summary vector). Panel B reports demeaned topic relevance scores for selected high- and low-relevance firms, color-coded by sector (Information sector in crimson, other sectors in slate grey). Both panels share identical aspect ratios ($1.5$).}\par}
\end{figure}

\subsubsection*{A.3.3. Research Outputs}

\begin{itemize}
    \item \textbf{10-K Item 1 Business Description Corpus:} Curated and verified annual report summaries.
    \item \textbf{Firm-Topic Relevance Panel:} Longitudinal exposure scores linking public firms to policy debates.
\end{itemize}

\newpage
\section*{Appendix B. Overview of Datasets Used for Data Pipeline}
\addcontentsline{toc}{section}{Appendix B. Overview of Datasets Used for Data Pipeline}

This appendix provides an integrated overview of the primary input corpora, verification benchmarks, and research output datasets utilized across our AI-enabled pipeline.

\subsection*{A. Primary Data Inputs}
\begin{itemize}
    \item \textbf{Congressional Hearings Dataset (U.S. GPO GovInfo CHRG Collection):} Complete package manifest of 36,378 hearings spanning 1995--2025. Focal sample covers 2,352 hearings before House Science and Senate Commerce committees containing 93.6 million words across 116,000 speaker turns.
    \item \textbf{Publicly Listed Firms Business Descriptions (SEC Form 10-K Filings):} Full SEC EDGAR archive accessed via sec-api.io, comprising 305,140 annual reports. Mapped to financial fundamentals through a bijective CRSP--Compustat crosswalk linking 25,207 CIK $\leftrightarrow$ GVKEY identifier pairs. Yields 22,299 S\&P 500 reports complete with extracted text and usability flags.
\end{itemize}

\subsection*{B. Verification Data}
\begin{itemize}
    \item \textbf{Independent Reference Abstracts (C-SPAN Archives API):} 18,697 committee program summaries covering 1995--2025 written by professional C-SPAN editors independently of congressional staff, serving as an external ground-truth benchmark for discriminant validity testing.
\end{itemize}

\subsection*{C. Output Research Datasets}
\begin{itemize}
    \item \textbf{Text Summarization Corpus (Stage 2 Outputs):} Cleaned, turn-attributed transcripts; intermediate chunk summaries; finalized hearing summaries (\texttt{hearing\_summaries.parquet}).
    \item \textbf{Policy Stance Panel Dataset (Stage 3 Outputs):} Standardized 111-topic policy taxonomy codebook (\texttt{frozen\_codebook\_111.parquet}); longitudinal member-topic stance panel (\texttt{member\_stance\_panel.parquet}).
    \item \textbf{Corporate Topic Relevance Dataset (Stage 4 Outputs):} Form 10-K Item 1 corporate business description embeddings; congressional debate centroid vectors; firm-year policy exposure panel (\texttt{firm\_topic\_relevance\_panel.parquet}).
\end{itemize}

\newpage
\section*{Appendix C. Example of Research Enabled by LLM Data Pipeline}
\addcontentsline{toc}{section}{Appendix C. Example of Research Enabled by LLM Data Pipeline}

\subsection*{Evaluating the Impact of Privacy Policy Debate Following the Cambridge Analytica Scandal on the Trajectory of Firm Corporate Research on Privacy Technology}

As described above, we demonstrate how the application of our data pipeline to congressional testimony and company 10-K filings allows us to quantify which firms were impacted by policy debate and study this question in a way not previously possible. 

Leveraging this pipeline, we generate an index providing both the stance of a policy debate (corporate and consumer data privacy) and the extent to which these policy positions impacted specific firms. Crucially, these issues would not have been possible to identify, let alone empirically study, in the absence of our staged approach: while prominent technology firms like Meta/Facebook were known targets following the Cambridge Analytica scandal, identifying the broader universe of affected firms across diverse sectors has historically been intractable.

In Figure 4, we evaluate patenting intensity across topic relevance quartiles ($Q_1$ to $Q_4$). Panel A plots the distribution of topic relevance scores against the intensity of privacy patenting. Firms identified as more exposed to the privacy debate generate an overwhelming share of privacy patents. Panel B evaluates annual patent filing rates before and after 2016 (2009--2015 vs. 2016--2020), showing that privacy innovation surged predominantly among Quartile 4 firms (rising from 5.57 to 8.84 patents per firm per year, $+3.27$ patents/year/firm), while lower quartiles remained flat. Panel C disaggregates these changes across specific technology subclasses, revealing surges in data anonymization (\texttt{G06F 21/6254}, \texttt{H04W 12/033}, $+335.1\%$) and network-level activity tracking (\texttt{H04L 67/535}, $+5,080\%$).

\begin{figure}[htbp]
    \centering
    \begin{minipage}{0.32\textwidth}
        \centering
        \textbf{\scriptsize Panel A: Distribution \& Patents}\par\vspace{2pt}
        \includegraphics[width=\textwidth]{figures/fig_app_c_panel_a_topic_relevance.png}
    \end{minipage}\hfill
    \begin{minipage}{0.32\textwidth}
        \centering
        \textbf{\scriptsize Panel B: Pre/Post 2016 by Quartile}\par\vspace{2pt}
        \includegraphics[width=\textwidth]{figures/fig_app_c_panel_b_quartiles.png}
    \end{minipage}\hfill
    \begin{minipage}{0.32\textwidth}
        \centering
        \textbf{\scriptsize Panel C: Pre/Post by Subclass}\par\vspace{2pt}
        \includegraphics[width=\textwidth]{figures/fig_app_c_panel_c_subclasses.png}
    \end{minipage}
    \caption{\textbf{Topic Relevance and Privacy Patenting Responses Pre- and Post-2016.}}
    \vspace{4pt}
    {\footnotesize\textit{Notes: Sample comprises 1,058 publicly traded firms scored for policy relevance. Panel A shows patent intensity by relevance quartile. Panel B shows pre- vs. post-2016 changes by quartile. Panel C shows annual filing changes across privacy technology subclasses.}\par}
\end{figure}

To evaluate the dynamic trajectory of this response, Figure 5 (referred to as Figure 6 in the pipeline) presents annual event-study estimates of patenting probability from 2009 to 2020 relative to 2015 across topic relevance quartiles ($Q_1$ to $Q_4$). The outcome in each panel is an indicator equal to one if firm $i$ filed at least one patent application in year $t$ in the focal technology group, estimated using firm fixed effects with standard errors clustered at the firm level.

\begin{figure}[htbp]
    \centering
    \begin{minipage}{0.48\textwidth}
        \centering
        \textbf{\footnotesize Panel A: Anonymizing / Decorrelating Data}\par\vspace{2pt}
        \includegraphics[width=\textwidth]{figures/fig_app_c_fig6_panel_a_anonymizing.png}
    \end{minipage}\hfill
    \begin{minipage}{0.48\textwidth}
        \centering
        \textbf{\footnotesize Panel B: Tracking User Activity}\par\vspace{2pt}
        \includegraphics[width=\textwidth]{figures/fig_app_c_fig6_panel_b_tracking.png}
    \end{minipage}
    \caption{\textbf{Annual Event-Study Estimates of Patenting Probability Relative to 2015 by Relevance Quartile.}}
    \vspace{4pt}
    {\footnotesize\textit{Notes: Sample comprises 1,058 firms. Estimated using firm fixed effects with standard errors clustered at the firm level. Vertical capless error bars represent 95\% confidence intervals. Vertical red dashed lines indicate the 2015/2016 policy transition and the 2018/2019 GDPR/CCPA enforcement milestone. Panel A evaluates Anonymizing or Decorrelating Data (CPC G06F 21/6254, H04W 12/033); Quartile 4 firms exhibit a sharp increase of $+6.0$ pp in 2019 ($p < 0.001$) and $+6.4$ pp in 2020 ($p < 0.001$). Panel B evaluates Tracking User Activity (CPC H04L 67/535); Quartile 4 firms experience a continuous rise up to $+9.8$ pp in 2020 ($p < 0.001$). Pre-trends are flat across all quartiles.}\par}
\end{figure}

In summary, this empirical exercise demonstrates how our AI-enabled pipeline constructs auditable, highly predictive firm-level measures that make it possible to study organizational innovation responses to policy debates.

\newpage
\section*{Appendix D. Reproducible and Adaptable Data Pipelines with Coding Agents}
\addcontentsline{toc}{section}{Appendix D. Reproducible and Adaptable Data Pipelines with Coding Agents}

While this data pipeline leverages AI tools to perform critical measurement tasks, an additional structural benefit is emerging: research pipelines shared via online repositories (like GitHub) in machine-readable formats can be rapidly deployed, verified, and adapted by autonomous coding agents.

Rather than treating papers as static documents, online repositories serve as active software environments where researchers share modular tools. To demonstrate this capability, the accompanying repository folder (\href{https://github.com/HC-Geopolitics-Innovation/LLM-Text-Analysis-and-Synthesis/tree/main/X2%20-%20Example%20of%20Reproducing%20and%20Adapting%20the%20Research%20Pipeline}{\texttt{X2 - Example of Reproducing and Adapting the Research Pipeline/}}) provides:
\begin{enumerate}
    \item \textbf{Sample Agent Prompt:} An actionable prompt (\href{https://github.com/HC-Geopolitics-Innovation/LLM-Text-Analysis-and-Synthesis/blob/main/X2%20-%20Example%20of%20Reproducing%20and%20Adapting%20the%20Research%20Pipeline/X1%20-%20Sample%20Agent%20Prompt%20For%20Adapting%20Research%20Pipeline%20To%20New%20Question.md}{\texttt{X1 - Sample Agent Prompt For Adapting Research Pipeline To New Question.md}}) instructing an LLM coding agent to read the repository architecture, inspect input corpora, and adapt the staged methodology to a new research setting (e.g., regulatory debate surrounding renewable energy and corporate carbon emission innovation).
    \item \textbf{Implementation Plan:} A full, researcher-evaluable implementation plan (\href{https://github.com/HC-Geopolitics-Innovation/LLM-Text-Analysis-and-Synthesis/blob/main/X2%20-%20Example%20of%20Reproducing%20and%20Adapting%20the%20Research%20Pipeline/X2%20-%20Sample%20Implementation%20Plan%20Generated%20by%20Pipeline.md}{\texttt{X2 - Sample Implementation Plan Generated by Pipeline.md}}) detailing step-by-step adaptations, schema crosswalks (from CPC privacy codes to Y02 climate mitigation patent classifications), verification checkpoints, and execution timelines.
\end{enumerate}

Researchers can provide this prompt to their own AI tools, evaluate the proposed implementation plan, and deploy adaptable pipelines for organizational and social science research.

\end{document}
